\section{Related Works}
\paragraph{Autoregressive VLA.}
Autoregressive VLA models(see~\Cref{fig:method_compare}) represent actions as discrete tokens, enabling LLM-style decoding and scalable training.
Representative systems include RT-1~\citep{brohan2022rt}, OpenVLA~\citep{kim2024openvla}, and LLaVA-VLA~\citep{song2026rethinking}, which unify vision, language, and action prediction under the next-to-next objective.
% Subsequent work improves perception or reasoning in discrete VLA pipelines, such as Robo-Flamingo~\citep{li2024vision}, GR-1~\citep{wu2023unleashing}, and 
ReconVLA~\citep{song2025reconvla} improves perception by reconstructing the current frame to assist discrete action generation.
Meanwhile, dynamic inference or efficiency tokenizer designs (e.g., Deer-VLA~\citep{yue2024deer} and FAST~\citep{pertsch2025fast}) target practical deployment constraints.
Methods based on vector quantization (VQ)~\citep{wang2025vq,liu2025faster,dong2026actioncodec} provide a flexible alternative by learning discrete latent representations.
UniVLA~\citep{univla} and WorldVLA~\citep{worldvla} discretize future frames to leverage dynamic visual information for action generation during both training and inference.
% 
These autoregressive models are effective but inherently sequential, which motivates our refinement alternative.
\paragraph{Discrete Diffusion VLA.}
Discrete diffusion VLA models extend dLLM-style~\citep{yu2025discrete} parallel denoising to action sequences, enabling full-attention, multi-token decoding with iterative refinement.
Representative methods differ in their corruption processes and efficiency strategies.
PD-VLA~\citep{song2025accelerating} follows a BART-style noising scheme by substituting tokens with vocabulary items and learning to reconstruct the sequence, while Discrete Diffusion VLA~\citep{liang2025discrete} and LLADA-VLA~\citep{wen2025llada} adopt BERT-style masking with a dedicated mask token~\citep{bert}.
To reduce inference cost, CEED-VLA~\citep{song2025ceed} applies consistency distillation to shrink the number of denoising steps with minimal performance loss.
On the data and modeling side, Dream-VLA~\citep{ye2025dream} scales discrete diffusion pretraining on the OXE~\citep{o2024open} datasets following OpenVLA~\citep{kim2024openvla}, and UD-VLA/dVLA~\citep{chen2025unified,wen2025dvla} incorporate visual or textual chain of thought (CoT) and jointly diffuse future frames, reasoning traces, and actions for unified perception--reasoning--action generation.
However, these dVLA methods cannot perform token-level refinement, as their updates are restricted to masked positions rather than allowing full-sequence refinement.
\paragraph{Discrete Flow Matching LLM and VLM.}
Discrete flow matching~\citep{gat2024discrete} provides a principled framework for modeling probability paths over discrete tokens, with theoretical work~\citep{shaul2024flow} establishing kinetic-optimal discrete paths and velocity formulations that generalize mask-based mixtures.
Building on these foundations, several large-scale multimodal models adopt discrete flow matching for unified understanding and generation, including Fudoki~\citep{wang2025fudoki} and Next-Omni~\citep{luo2025next}, demonstrating strong any-to-any or omnimodal capabilities.
EditFlow~\citep{havasi2025edit} introduces explicit edit operations (insertion, replacement, deletion) with learned rates to enable flexible sequence refinement, while OneFlow~\citep{nguyen2025oneflow} extends this idea to concurrent mixed-modal and interleaved generation.
URSA~\citep{deng2025uniform} further illustrates how discrete metric structures can guide token transitions in generative modeling.
We introduce discrete flow matching to the action modality, leveraging iterative refinement via a velocity field to progressively enhance action generation.